% !TEX TS-program = xelatex
% !BIB TS-program = bibtex
\documentclass[12pt,letterpaper]{article}
\usepackage{style/qss20reportstyle} % import dsc180reportstyle.sty

% Grey italic [WRITE] prompts: what each paragraph should say, with the
% numbers to use. DELETE every [WRITE] block before submitting.
\newcommand{\writeprompt}[1]{{\color{gray}\textit{[WRITE] #1}}}
\newcommand{\authnote}[1]{{\color{gray}\textit{[NOTE] #1}}}

%%%%%%%%%%%%%%%%%%%%%%%%%%%%%%%%%%%%%%%%%%%%%%%%%%%%%%%%
%%%% Title and Authors
%%%%%%%%%%%%%%%%%%%%%%%%%%%%%%%%%%%%%%%%%%%%%%%%%%%%%%%%

%%% New title: the old one ("The Spread Doesn't Tank") overclaims.
%%% The spread DID is a null in the main spec and insignificant in the
%%% named-five spec, so the title is now a question, not a claim.
\title{Tanking and the Spread: Do Betting Markets Price Late-Season Losing?}

\author{Sid Singh \\
  {\tt siddharth.v.singh.27@dartmouth.edu} \\\And
  Ahaan Jindal \\
  {\tt ahaan.s.jindal.27@dartmouth.edu}}

\begin{document}
\maketitle

%%%%%%%%%%%%%%%%%%%%%%%%%%%%%%%%%%%%%%%%%%%%%%%%%%%%%%%%
%%%% Abstract and Links
%%%%%%%%%%%%%%%%%%%%%%%%%%%%%%%%%%%%%%%%%%%%%%%%%%%%%%%%

\vspace{-1cm}
\begin{center}
GitHub Code Repository: \url{https://github.com/asj005/QSS20-Project}
\end{center}

\maketoc
% \clearpage

%%%%%%%%%%%%%%%%%%%%%%%%%%%%%%%%%%%%%%%%%%%%%%%%%%%%%%%%
%%%% Main Contents
%%%%%%%%%%%%%%%%%%%%%%%%%%%%%%%%%%%%%%%%%%%%%%%%%%%%%%%%


\section{Introduction}

Under the NBA's draft lottery system, the teams with the worst records received the best
odds at the top pick. This therefore creates a strong incentive system that rewards losing
late in the season (a phenomenon called ``tanking''). The 2025--26 season was an extreme
case: five teams (Washington, Indiana, Memphis, Brooklyn and Utah) won fewer than 18\%
of their games after the All-Star break (Associated Press 2026). In May 2026, to avoid this
issue and preserve the spirit of the game, the league's Board of Governors approved a new
``3-2-1'' lottery that cuts the odds of the three worst teams, beginning with the 2027 draft
(Associated Press 2026).
Betting markets give us a way to measure tanking that does not rely on accusations or
anecdotes. The point spread and the total line are the market's forecasts of each game's
margin and combined score. If markets fully anticipated tanking, tanking teams would still
beat the spread about half the time. Our research questions are:
1. After the All-Star break, do teams in tanking position underperform the point spread
and the total line, relative to their own pre-break performance and to other teams?
2. If so, how is the mispricing produced? Do these teams stop scoring, stop defending,
or rest their stars?
3. Does a crowd-driven prediction market (Polymarket) price tanking any better than
sportsbooks?
Our method is a difference-in-differences design on market forecast errors, with treatment
status defined from information available before each game.

\writeprompt{Closing paragraph: preview the answers in one sentence each, then a
roadmap sentence. The spread answer is a null in the main spec (DID $-0.04$
points, $t=-0.03$); the totals answer leans toward a miss ($+3.28$ points,
$t=1.37$, not significant); the mechanism is defense, not scoring or rest
(points allowed DID $+3.00$, $t=2.16$); Polymarket is no better than the
sportsbook on tanking games (Brier 0.1190 vs 0.1198 post-break). One sentence
per section for the roadmap.}

\section{Related Work}

\authnote{This section cost 0.5 points at Checkpoint 2 for bold headers and
bullet points. Write it as continuous prose: one paragraph per paper or per
strand, and blend how each paper differs from yours into the description of
that paper instead of listing differences separately. No bold subheadings, no
bullets. Reuse the papers from Checkpoint 2; no citations are invented here.}

\writeprompt{Paragraph 1: the tanking literature. What is known about
tanking as a strategy: lottery incentives, how teams tank (rest, rotations,
effort), and the league's responses, including the 2019 odds flattening and
the 2026 ``3-2-1'' reform. Your paper's contribution is measurement: market
forecast errors instead of accusations.}

\writeprompt{Paragraph 2: sports-betting market efficiency. Whether
sportsbook lines are efficient and where inefficiencies have been claimed
(favorite-longshot bias, home bias, late-season effects). Your spread and
totals DID speak directly to this: a significant post-break deterioration
for tanking teams would be a documented inefficiency; a null is evidence
the market adjusted.}

\writeprompt{Paragraph 3: prediction markets vs sportsbooks. What the
literature finds on whether crowd-driven prices beat bookmaker prices, and
in which settings. Your RQ3 tests that head-to-head on the same games,
scored with Brier scores, split by tanking status and by pre/post break.
End the section with a sentence stating the gap your paper fills.}

\section{Data}

\writeprompt{Paragraph 1: the panel in one paragraph. 1,230 games, the full
2025--26 NBA regular season, October 21, 2025 to April 12, 2026, 30 teams,
82 games each, one row per game. Scores, closing moneylines, spreads and
totals; Polymarket moneyline, spread and totals prices taken as the last
trade before tip-off; box scores and officials for every game. The NBA Cup
final (December 16, 2025) is excluded because it does not count in the
regular-season standings; playoffs and play-in games are not in the sample.
The All-Star break falls between February 12 and February 19, 2026 (seven
days with no games); pre-break is 819 games, post-break 411.}

Data sources (hyperlinked):
\begin{itemize}
\item Scores, box scores, officials, tip-off times: \url{https://www.basketball-reference.com/leagues/NBA_2026_games.html}
\item Closing moneylines, spreads, totals: \url{https://www.teamrankings.com/nba/matchup/}
\item Prediction-market prices: Polymarket Gamma API (\url{https://gamma-api.polymarket.com/events?tag_id=745}, NBA tag 745) and the CLOB price-history endpoint, last price at or before tip-off.
\end{itemize}

\authnote{Why Polymarket and not Kalshi: Kalshi publishes no public price
history for event markets, so historical Kalshi odds for the 2025--26 season
do not exist. Say so in one sentence; it pre-empts the obvious question.}

\writeprompt{Paragraph 2: what each source contributes and how they join.
Everything joins on date and the two teams; team names match exactly across
sources. Moneyline coverage is nearly complete (Table 2); spread and totals
coverage on Polymarket is thinner because those markets sometimes had no
trade in the 24 hours before tip-off. Implied probabilities are vig-free:
each moneyline is converted to an implied probability and the two sides are
normalized to sum to one (average overround 4.50\%).}

\begin{table}[htbp]
\centering
\caption{Summary statistics for key variables (1,230 games).}
\label{tab:summary}
\begin{tabular}{lrrrr}
\hline
Variable & $n$ & Mean & Std.\ dev.\ & Range \\
\hline
Home margin (points) & 1,230 & 1.74 & 16.41 & $-$55 to 54 \\
Actual total (points) & 1,230 & 231.22 & 19.88 & 169 to 302 \\
Sportsbook total line & 1,230 & 230.73 & 7.67 & 206.5 to 254.5 \\
Sportsbook implied home-win prob.\ & 1,230 & 0.550 & 0.219 & 0.07 to 0.95 \\
Polymarket implied home-win prob.\ & 1,229 & 0.552 & 0.225 & 0.05 to 0.98 \\
Total fouls per game & 1,230 & 39.73 & 6.85 & 14 to 70 \\
Home win rate & 1,230 & 0.554 & --- & --- \\
Over rate (actual $>$ line) & 1,230 & 0.491 & --- & --- \\
Home ATS cover rate & 1,230 & 0.516 & --- & --- \\
\hline
\end{tabular}
\end{table}

\begin{table}[htbp]
\centering
\caption{Coverage: non-missing observations per variable.}
\label{tab:coverage}
\begin{tabular}{lrrr}
\hline
Variable & $n$ & Games & Share \\
\hline
Sportsbook moneyline (all games) & 1,230 & 1,230 & 1.000 \\
Polymarket moneyline & 1,229 & 1,230 & 0.999 \\
Polymarket spread line & 1,221 & 1,230 & 0.993 \\
Polymarket spread prob.\ & 1,119 & 1,230 & 0.910 \\
Polymarket total line & 1,221 & 1,230 & 0.993 \\
Polymarket over prob.\ & 1,063 & 1,230 & 0.864 \\
Both top scorers' trailing form & 646 & 1,230 & 0.525 \\
Form-test sample (form + no star DNP) & 581 & 1,230 & 0.472 \\
\hline
\end{tabular}
\end{table}

\writeprompt{Paragraph 3: read Tables 1--2 in words. Home teams won 55.4\%
of games by an average of 1.74 points; 49.1\% of games went over the total
and home teams covered 51.6\% of spreads, so the lines look roughly right on
average. The average total line was 230.7 points against an average actual
total of 231.2. Polymarket moneyline coverage is 1,229 of 1,230 games; the
one miss is March 1, 2026, San Antonio at New York, which had no trade in
the 24 hours before tip-off.}

\begin{figure}[htbp]
\centering
\includegraphics[width=\linewidth]{figure/fig6_eda.png}
\caption{Exploratory data analysis. (a) Home margin centers near zero with
mean $+1.74$. (b) Actual totals track the sportsbook line. (c) The two
markets' implied-probability distributions nearly coincide. (d) Total fouls
per game average 39.7.}
\label{fig:eda}
\end{figure}

\writeprompt{Paragraph 4: EDA, before any test (this was a missing piece at
Checkpoint 2). Walk through Figure 1: the margin distribution, the tight
relationship between the total line and the actual total, and the
near-identical implied-probability distributions, which already hints that
the two markets will score similarly. State the tanking fact that motivates
the design: the five named teams won 10--17\% of their games after the
break (Washington 3--26, Indiana 4--23, Memphis 5--24, Brooklyn 5--24, Utah
4--22), verified from the game data.}

\authnote{Cleaning detail lives in Appendix A, not here. That was a
0.3-point deduction last time. Do not move it back.}


\section{Methodology}

\authnote{This section cost 0.5 points at Checkpoint 2: methods were stated,
not explained formally, and there was no EDA before the tests. Every measure
below is defined with its formula, and every regression is written out.
Describe the method behind each result in prose, not just the result.}

\subsection{Treatment: who is tanking}

\writeprompt{One paragraph, and it matters: get the two definitions straight.
The MAIN treatment group is the five teams with the worst PRE-BREAK win
percentage: Sacramento (21.4\%), Washington (26.4\%), New Orleans (26.8\%),
Indiana (27.3\%), Brooklyn (28.3\%). Treatment is therefore defined from
information available before each game, as the Introduction promises. The
five teams the Introduction names (Washington, Indiana, Memphis, Brooklyn,
Utah) are the ROBUSTNESS group: they are selected on POST-break win
percentage (all below 18\%), so their selection uses post-break outcomes.
Memphis (37.7\% pre-break) and Utah (32.1\%) were not bottom-five at the
break; Sacramento and New Orleans were. Both definitions are reported
throughout, and the results agree in direction.}

\subsection{Outcomes: market forecast errors}

\writeprompt{One paragraph defining the two forecast errors. Let $s_g$ be the
home team's spread in game $g$ (negative when the home team is favored) and
$m_g$ the home margin in points. The home team's ATS (against-the-spread)
margin is $m_g + s_g$; the away team's is its negative. A team covers when
its ATS margin is positive (4 of 2,460 team-games push exactly). The total
forecast error is the actual combined score minus the total line. Both are
defined so that a positive value means the market UNDER-estimated that side.}

\begin{align}
\text{ATS margin}_{ig} &= m_g + s_g \quad \text{(home team)} \label{eq:ats} \\
\text{Total error}_{g} &= \text{actual total}_g - \text{total line}_g \label{eq:toterr}
\end{align}

\subsection{The difference-in-differences design}

\writeprompt{One paragraph. The unit is the team-game (2,460 rows: two per
game). The DID compares tanking teams to other teams, after the break to
before it. The interaction coefficient $\delta$ is the estimate: how much
worse (or better) tanking teams did against the market after the break,
relative to their own pre-break performance and to other teams' change over
the same window. Report the specification with team fixed effects as the
main one; the tanking main effect is absorbed by the team effects there.}

\begin{equation}
Y_{it} = \alpha + \beta \cdot \text{Tank}_i + \gamma \cdot \text{Post}_t + \delta \cdot (\text{Tank}_i \times \text{Post}_t) + \varepsilon_{it} \label{eq:did}
\end{equation}
\begin{equation}
Y_{it} = \alpha_i + \gamma \cdot \text{Post}_t + \delta \cdot (\text{Tank}_i \times \text{Post}_t) + \theta \cdot \text{Home}_{it} + \varepsilon_{it} \label{eq:didfe}
\end{equation}

\writeprompt{One paragraph on standard errors. Each game contributes two
team-game rows whose ATS margins sum to zero by construction, so standard
errors are clustered by game (1,230 clusters). The totals regressions are
run at the game level (1,230 observations) and need no clustering. State
this explicitly; it is the kind of detail the grader looks for.}

\writeprompt{One paragraph on the identifying assumption. The DID assumes
tanking and non-tanking teams would have trended together without the break
(parallel trends). Figure 2 checks it: monthly means for ATS margin and
points allowed, before and after the break. Report the check honestly. For
points allowed, the outcome with the significant DID, there is no
differential pre-trend (joint F-test on the pre-break interactions: F = 1.32,
p = 0.26). For ATS margin the test is marginal (F = 2.33, p = 0.055),
driven by December, when tanking teams beat the spread by 2.4 points on 65
team-games; monthly samples are small (25--80 team-games), so the monthly
means are noisy, and with four pre-period interactions one t-statistic near
2.5 is roughly what chance produces. State both results and the December
caveat rather than hiding the marginal one.}

\begin{figure}[htbp]
\centering
\includegraphics[width=\linewidth]{figure/tank_fig5_pretrends.png}
\caption{Monthly means by tanking status (main spec), before and after the
All-Star break. (a) ATS margin. (b) Points allowed per game. The dashed
line marks the break.}
\label{fig:pretrends}
\end{figure}

\subsection{Mechanisms and the prediction-market comparison}

\writeprompt{One paragraph on mechanisms. The same DID is run with three
different outcomes: points scored, points allowed, and an indicator for
whether one of the team's top-three scorers (by season points per game) sat
out. Scoring, defense, and rest are the three channels the Introduction
names, and the three DIDs separate them.}

\writeprompt{One paragraph on RQ3. Polymarket and sportsbook moneyline
prices are scored with the Brier score, split four ways: tanking-team games
vs other games, pre-break vs post-break. If the crowd prices tanking better,
its Brier score should be lower on tanking games after the break, relative
to the sportsbook. The spread lines are also compared directly (mean
absolute difference 0.25 points, correlation 0.992).}

\begin{equation}
\text{Brier} = \frac{1}{N} \sum_{i} (p_i - y_i)^2 \label{eq:brier}
\end{equation}

\section{Results \& Findings}

\subsection{The spread: no tanking blind spot}

\begin{figure}[htbp]
\centering
\includegraphics[width=\linewidth]{figure/tank_fig1_spread.png}
\caption{Spread outcomes by tanking status (main spec: worst five pre-break
win\%) and period. (a) Cover rate. (b) ATS margin in points. Tanking teams
cover no worse after the break.}
\label{fig:spread}
\end{figure}

\begin{table}[htbp]
\centering
\caption{Spread DID cells, main spec (worst five by pre-break win\%).}
\label{tab:spread}
\begin{tabular}{llrrrr}
\hline
& & $n$ & Cover rate & ATS margin & Win rate \\
\hline
Other teams & Pre-break & 1,365 & 0.506 & 0.24 & 0.548 \\
& Post-break & 685 & 0.505 & 0.25 & 0.552 \\
Tanking teams & Pre-break & 273 & 0.465 & $-$1.20 & 0.260 \\
& Post-break & 137 & 0.474 & $-$1.24 & 0.241 \\
\hline
\end{tabular}
\end{table}

\writeprompt{Paragraph 1: read Figure 3 and Table 3. Tanking teams covered
46.5\% pre-break and 47.4\% post-break; other teams covered 50.5\% in both
periods. The ATS margin for tanking teams is about $-$1.2 points in BOTH
periods, against $+$0.24 for other teams. The DID estimate is $-0.04$ points
($t=-0.03$, $p=0.98$) with team fixed effects, and $+0.01$ on the cover
indicator ($t=0.21$). This is a clean null: the spread did not get worse for
tanking teams after the break. Note the level effect honestly: tanking teams
underperform the spread by about 1.2 points all season, in both periods, so
the market shades them a little light throughout; what does NOT happen is a
post-break deterioration.}

\writeprompt{Paragraph 2: the named-five robustness check (Appendix Table
A1). There the direction flips toward the title's story: cover rate 48.5\%
pre-break to 41.4\% post-break, ATS margin $-0.56$ to $-1.85$, DID $-1.55$
points ($t=-1.09$, $p=0.27$) and $-0.085$ on the cover indicator ($t=-1.47$,
$p=0.14$). Report it as suggestive but not significant: with 140 post-break
team-games, the standard error is about 1.4 points, so the data cannot rule
out a zero effect. Do not oversell this; the main spec is a zero and the
robustness spec is a noisy negative.}

\subsection{The totals: a suggestive miss}

\begin{figure}[htbp]
\centering
\includegraphics[width=\linewidth]{figure/tank_fig2_totals.png}
\caption{Totals outcomes by whether the game involves a tanking team (main
spec) and period. (a) Over-hit rate. (b) Total forecast error in points.
Tanking games go over more after the break.}
\label{fig:totals}
\end{figure}

\begin{table}[htbp]
\centering
\caption{Totals DID cells, main spec (games involving a tanking team).}
\label{tab:totals}
\begin{tabular}{llrrrr}
\hline
& & $n$ & Over-hit & Total error & Actual total \\
\hline
Other games & Pre-break & 562 & 0.479 & 0.01 & 231.13 \\
& Post-break & 284 & 0.486 & 0.50 & 230.05 \\
Tanking games & Pre-break & 257 & 0.483 & $-$0.07 & 231.04 \\
& Post-break & 127 & 0.575 & 3.70 & 234.55 \\
\hline
\end{tabular}
\end{table}

\writeprompt{One or two paragraphs. Games involving a tanking team went over
57.5\% of the time after the break, against 48.3\% before it and 48.6\% for
other games after the break. The total forecast error for tanking games rose
from $-0.07$ to $+3.70$ points. The DID is $+3.28$ points ($t=1.37$,
$p=0.17$) and $+0.085$ on the over-hit indicator ($t=1.30$, $p=0.19$). Same
direction in the named-five spec ($+3.11$ points, $t=1.30$; over-hit
$+0.106$, $t=1.63$, $p=0.10$). This is the closest thing in the paper to a
market miss, but it is not significant at conventional levels: say so
plainly, and point to the mechanism section for why the direction makes
sense anyway (worse defense means more points).}

\subsection{The mechanism is defense}

\begin{figure}[htbp]
\centering
\includegraphics[width=\linewidth]{figure/tank_fig3_mechanism.png}
\caption{Mechanism outcomes by tanking status (main spec) and period. (a)
Points allowed per game. (b) Share of team-games with a top-three scorer
sitting out.}
\label{fig:mechanism}
\end{figure}

\begin{table}[htbp]
\centering
\caption{Mechanism DID cells, main spec (team-game level).}
\label{tab:mechanism}
\begin{tabular}{llrrrrr}
\hline
& & $n$ & Scored & Allowed & Net & Star DNP \\
\hline
Other teams & Pre-break & 1,365 & 116.43 & 114.75 & 1.68 & 0.081 \\
& Post-break & 685 & 116.40 & 114.42 & 1.98 & 0.114 \\
Tanking teams & Pre-break & 273 & 111.17 & 119.55 & $-$8.39 & 0.066 \\
& Post-break & 137 & 112.34 & 122.23 & $-$9.89 & 0.131 \\
\hline
\end{tabular}
\end{table}

\begin{table}[htbp]
\centering
\caption{DID estimates (interaction coefficient, $t$-statistic in
parentheses). Main spec: worst five by pre-break win\%. Named spec: the
Introduction's five teams. Standard errors clustered by game in the
team-game regressions. $^{*}$ $p<0.10$, $^{**}$ $p<0.05$, $^{***}$ $p<0.01$.}
\label{tab:did}
\begin{tabular}{lrr}
\hline
Outcome & Main spec & Named-five spec \\
\hline
ATS margin (points) & $-$0.04 ($-0.03$) & $-$1.55 ($-1.09$) \\
Cover indicator & 0.01 (0.21) & $-$0.08 ($-1.47$) \\
Total forecast error (points) & 3.28 (1.37) & 3.11 (1.30) \\
Over-hit indicator & 0.09 (1.30) & 0.11 (1.63$^{*}$) \\
Points scored & 1.20 (0.85) & $-$1.08 ($-0.80$) \\
Points allowed & 3.00 (2.16$^{**}$) & 5.17 (3.66$^{***}$) \\
Star DNP rate & 0.03 (0.95) & 0.07 (1.72$^{*}$) \\
\hline
\end{tabular}
\end{table}

\writeprompt{Paragraph 1: the three channels, separated. Scoring: the DID is
$+1.20$ points ($t=0.85$), and in the named spec $-1.08$ ($t=-0.80$):
tanking teams did NOT stop scoring. Defense: the DID is $+3.00$ points
allowed ($t=2.16$, $p=0.031$), and $+5.17$ ($t=3.66$, $p<0.001$) in the named
spec: this is the significant result in the paper, in both definitions.
Rest: the star DNP rate rose for everyone after the break (other teams 8.1\%
to 11.4\%), and the tanking-specific increase is $+3.3$ points ($t=0.95$),
$+6.7$ points ($t=1.72$, $p=0.085$) in the named spec. So the mechanism is
defense first, rest maybe, scoring no. Their net rating worsened from
$-8.39$ to $-9.89$ (main spec) and from $-7.12$ to $-12.26$ (named spec).}

\writeprompt{Paragraph 2: tie the mechanism back to the totals result. If
tanking teams defend worse but score the same, games involving them should
go over more often, which is exactly the direction of the totals DID. The
mechanism and the totals result corroborate each other even though the
totals DID alone is not significant.}

\subsection{Polymarket is no better on tanking games}

\begin{figure}[htbp]
\centering
\includegraphics[width=\linewidth]{figure/tank_fig4_brier.png}
\caption{Brier scores (lower is better) for Polymarket vs sportsbook
moneylines, split by whether the game involves a tanking team (main spec)
and period.}
\label{fig:brier}
\end{figure}

\begin{table}[htbp]
\centering
\caption{Brier scores, Polymarket vs sportsbook, by tanking status and period (main spec).}
\label{tab:brier}
\begin{tabular}{llrrr}
\hline
& & $n$ & Polymarket & Sportsbook \\
\hline
Tanking games & Pre-break & 257 & 0.1776 & 0.1780 \\
& Post-break & 127 & 0.1190 & 0.1198 \\
Other games & Pre-break & 562 & 0.2314 & 0.2309 \\
& Post-break & 284 & 0.1723 & 0.1730 \\
\hline
\end{tabular}
\end{table}

\writeprompt{One paragraph. The Brier scores are nearly identical in all
four cells: on tanking games after the break, Polymarket 0.1190 against the
sportsbook's 0.1198. On the spread, Polymarket's implied home-cover
probability for post-break tanking games averages 0.497 against an actual
home-cover rate of 0.527 ($n=110$), and the two markets' spread lines differ
by 0.25 points on average (correlation 0.992). The crowd does not see
tanking coming any better than the bookmaker does. Note why Brier scores
are lower on tanking games in general: heavy underdogs mostly lose as
expected, so the outcome is easier to predict.}

\section{Discussions \& Interpretations}

\writeprompt{Paragraph 1: what the spread null means. The cleanest reading
is that the market knew: tanking teams were already underperforming the
spread by about 1.2 points before the break, and the break changed nothing.
The spread had already priced in that these teams are bad. That is the
efficient-market result, and it is the paper's answer to RQ1: no, tanking
teams do not underperform the spread after the break relative to before it.
The named-five spec leans the other way but is not significant; report both
and let the reader see the uncertainty.}

\writeprompt{Paragraph 2: the defense mechanism and the totals suggestion.
Tanking in 2025--26 looks like not defending, not not scoring and not
resting stars much more than anyone else. That is consistent with how
tanking actually works in practice (young players getting minutes, effort
dropping on the defensive end) and it explains the totals direction. The
totals market looks like the one place the price may lag the behavior, but
the evidence is suggestive, not conclusive.}

\writeprompt{Paragraph 3: limitations. One season, so everything is a
2025--26 statement. The post-break tanking sample is small (127--137
games), which is why the totals and named-spec spread estimates are noisy.
The two tanking definitions disagree on Memphis, Utah, Sacramento and New
Orleans; the results agree in direction, but the treatment definition is a
judgment call and should be stated as one. Standard errors are clustered by
game but the design is still observational: the DID assumes tanking and
non-tanking teams would have trended together without the break (parallel
trends), and nothing in one season can prove that. The form tests from the
earlier analysis use 581 games for the coverage reasons in Section 3.}

\section{Conclusions \& Future Work}

\writeprompt{Paragraph 1: the findings in four sentences: the spread null,
the totals suggestion, the defense mechanism, the Polymarket tie. No new
numbers here; point back to the tables and figures.}

\writeprompt{Paragraph 2: future work. Three concrete directions. (1) More
seasons, to check whether the totals miss and the named-spec spread gap are
real or noise, and whether the 2027 ``3-2-1'' lottery changes tanking
behavior. (2) Player-level data: minutes and lineups for young players on
tanking teams after the break, to pin down the defense mechanism. (3)
In-game prices: whether live markets move fast enough when a tanking team's
defense collapses mid-game. Each follows from a specific limitation in
Section 6.}


%%%%%%%%%%%%%%%%%%%%%%%%%%%%%%%%%%%%%%%%%%%%%%%%%%%%%%%%
%%%% Reference / Bibliography
%%%%%%%%%%%%%%%%%%%%%%%%%%%%%%%%%%%%%%%%%%%%%%%%%%%%%%%%

\makereference
\bibliography{reference}
\bibliographystyle{style/qss20bibstyle}

\authnote{Add to reference.bib: the two Associated Press (2026) pieces from
Checkpoint 2, plus the data sources. BibTeX entries for the data sources
are in report/reference\_bib\_additions.bib in the repo. The lottery claim
is verified: the Board of Governors approved the ``3-2-1'' system on May 28,
2026, 29--1 (Memphis dissenting), effective for the 2027 draft.}

%%%%%%%%%%%%%%%%%%%%%%%%%%%%%%%%%%%%%%%%%%%%%%%%%%%%%%%%
%%%% Appendix
%%%%%%%%%%%%%%%%%%%%%%%%%%%%%%%%%%%%%%%%%%%%%%%%%%%%%%%%


% Optional, comment this part out if not needed

\clearpage
\makeappendix

\subsection{Named-five robustness tables}

\begin{table}[htbp]
\centering
\caption{Spread DID cells, named-five spec (Washington, Indiana, Memphis,
Brooklyn, Utah).}
\label{tab:spread-named}
\begin{tabular}{llrrrr}
\hline
& & $n$ & Cover rate & ATS margin & Win rate \\
\hline
Other teams & Pre-break & 1,368 & 0.502 & 0.11 & 0.539 \\
& Post-break & 682 & 0.518 & 0.38 & 0.572 \\
Tanking teams & Pre-break & 270 & 0.485 & $-$0.56 & 0.304 \\
& Post-break & 140 & 0.414 & $-$1.85 & 0.150 \\
\hline
\end{tabular}
\end{table}

\begin{table}[htbp]
\centering
\caption{Totals DID cells, named-five spec.}
\label{tab:totals-named}
\begin{tabular}{llrrrr}
\hline
& & $n$ & Over-hit & Total error & Actual total \\
\hline
Other games & Pre-break & 566 & 0.474 & $-$0.46 & 230.21 \\
& Post-break & 281 & 0.473 & 0.05 & 228.94 \\
Tanking games & Pre-break & 253 & 0.494 & 0.98 & 233.09 \\
& Post-break & 130 & 0.600 & 4.60 & 236.85 \\
\hline
\end{tabular}
\end{table}

\begin{table}[htbp]
\centering
\caption{Mechanism DID cells, named-five spec (team-game level).}
\label{tab:mechanism-named}
\begin{tabular}{llrrrrr}
\hline
& & $n$ & Scored & Allowed & Net & Star DNP \\
\hline
Other teams & Pre-break & 1,368 & 116.07 & 114.66 & 1.41 & 0.079 \\
& Post-break & 682 & 116.44 & 113.92 & 2.52 & 0.106 \\
Tanking teams & Pre-break & 270 & 112.94 & 120.07 & $-$7.12 & 0.078 \\
& Post-break & 140 & 112.24 & 124.50 & $-$12.26 & 0.171 \\
\hline
\end{tabular}
\end{table}

\subsection{Data cleaning and validation}

Cleaning detail lives here, not in the main text. All four sources join on
date and the two team names, which match exactly across Basketball-Reference,
TeamRankings, and Polymarket. Tip-off times come from the
Basketball-Reference schedule pages. Polymarket events come from the Gamma
API (NBA tag 745); prices come from the CLOB price-history endpoint as the
last traded price at or before tip-off (median 1.0 hour before tip-off; 1.3\%
of prices more than 6 hours stale). The Gamma fields \texttt{clobTokenIds}
and \texttt{outcomes} are JSON-encoded strings and must be parsed before use.
Player box scores cover 32,166 player-games, including 5,515 DNP rows kept
with \texttt{mp=`DNP'} and blank stats. A top scorer with no box-score row at
all (out injured or resting, not even listed) has missing trailing form; that
is the main reason the form sample is 581 games rather than 1,230. The DNP
flag fires only when the player has a row marked DNP. Validation: team points
from the box scores match the scores file in all 1,230 games; player points
sum to team totals in all 2,460 team-games. Five randomly chosen scores and
five randomly chosen odds lines were re-checked live against the sources.
Polymarket and sportsbook implied probabilities correlate at 0.997 with a
mean absolute difference of 0.012, which cross-validates the two price series
against each other. Known quirk, kept as listed: February 11, 2026, New York
at Philadelphia, where TeamRankings' moneyline (Philadelphia $-140$) and
spread (New York $-2.5$) disagree about the favorite. New York won 138--89.

\subsection{Repository map}

Notebooks 01--04 build the data end to end and reproduce every number in
this report: 01 odds and scores, 02 box scores, 03 Polymarket prices, 04 the
merge and all analysis. \texttt{tanking\_analysis.py} runs the DID design in
Section 4 under both tanking definitions and writes
\texttt{results/tanking\_results.txt}. \texttt{make\_figures.py} regenerates
Figures 1--8 and Tables 1--7 from the merged file and asserts the headline
numbers against \texttt{results/analysis\_results.txt} before writing
anything; \texttt{make\_tanking\_figures.py} regenerates the tanking
figures (Figures 3--6) and \texttt{pretrends\_check.py} regenerates the
pre-trends figure (Figure 2) and writes
\texttt{results/pretrends\_results.txt}. Raw
files live in \texttt{data/raw/}, the merged panel in \texttt{outputs/},
figures in \texttt{figures/}, and tables in \texttt{results/tables/}.

\end{document}
